Ath. að eðlisfræðingarnir þurfa að sitja sem einn hópur af 5 eða hópur af 2 og hópur af 3. 

\leftskip 15pt
\noindent Ef einn hópur af fimm þá sitja allar brautir saman og eðlisfræðingarnir eru á milli stærðfræði og hagnýt stærðfræði, 2 tilvik
\newline Ef hópur af 2 og 3.

\leftskip 30pt
\noindent G.r.f. að stærðfræðingarnir sitja fyrr á bekknum. Þá alveg fyrst get það verið hópur frá 0 til 6 stærðfræðingum,

\leftskip 45pt
\noindent Ef hópur 1 til 5 fremst þá eru tveir möguleikar fyrir hvar hinn hópur af stærðfræði sitjast, 10 tilvik
\newline Ef hópur 0 fremst þá þurfa stærðfræði að koma næst þ.a. 1 tilvik
\newline Ef hópur 6 fremst þá get hagnýtu stærðfræðingarnir skipt sér á 11 vegu þar sem einhverjir þurfa að vera í miðjunni þ.a. 11 tilvik

\leftskip 30pt
\noindent
22 tilvik ef stærðfræði er á undan hag
\newline G.r.f. að hagnýtu stærðfræðingarnir sitja fyrr á bekknum. Þá alveg fyrst get það verið hópur frá 0 til 11 hagnýtum stærðfræðingum,

\leftskip 45pt
\noindent Ef hópur 1 til 10 fremst þá eru tveir möguleikar fyrir hvar hinn hópur af stærðfræði sitjast, 20 tilvik
\newline Ef hópur 0 fremst þá þurfa hagnýt stærðfræði að koma næst þ.a. 1 tilvik
\newline Ef hópur 11 fremst þá get stærðfræðingarnir skipt sér á 6 vegu þ.a. 6 tilvik

\leftskip 30pt \noindent 27 tilvik ef hagnýt stærðfræði er á undan stæ.

\leftskip 15pt \noindent 49 tilvik saman og þá sinnum 2 þar sem hægt er að hafa 2 eða eðlisfræðinga fyrst þ.a. 98 tilvik

\leftskip 0pt\noindent Þá höfum við 100 leiðir til að raða eftir braut þ.a. $100\cdot 6!\dot11!\cdot5!$ leiðir til að raða nemendunum.

